\section{Summary}

In summary, we report the highest statistics measurement to date of neutrino neutral current single pion production on argon, including the first exclusive measurements of this process ever made in argon. These cross sections are measured using the MicroBooNE detector exposed to the Fermilab Booster Neutrino Beamline, which has $\langle E_{\nu} \rangle <1$ GeV. As presented within this paper,  kinematic distributions of the $\pi^0$ momentum and angle relative to the beam direction provide some sensitivity to contributions to this process from coherent and non-coherent pion production and suggest that, given the currently analyzed MicroBooNE data statistics, the nominal \textsc{genie} neutrino event generator used for MicroBooNE Monte Carlo modeling describes the observed distributions within uncertainties. This has provided an important validation check justifying the use of this sample as a powerful constraint for backgrounds to single-photon searches in MicroBooNE, \textit{e.g.}~in \cite{MicroBooNE:2021zai}. 

Using a total of 1,130 observed NC $\pi^0$ events, a flux-averaged cross section has been extracted for neutrinos with a mean energy of 804 MeV and has been found to correspond to $1.243\pm0.185$ (syst) $\pm0.076$ (stat), $0.444\pm0.098\pm0.047$, and $0.624\pm0.131\pm0.075$ $[10^{-38}\textrm{cm}^2/\textrm{Ar}]$ for the semi-inclusive NC$\pi^0$, exclusive NC$\pi^0$+1p, and exclusive NC$\pi^0$+0p processes compared to $1.678$, $0.722$, and $0.774$ $[10^{-38}\textrm{cm}^2/\textrm{Ar}]$ in the default \textsc{genie} prediction used by MicroBooNE. Comparison to other generators including \textsc{neut} and \textsc{NuWro} show reasonable agreement with the \textsc{neut} predictions found to be slightly more consistent with the MicroBooNE data-extracted cross-section for all three exclusive and semi-inclusive processes.

